\section{Arquitectura del sistema}
\label{sec:arquitectura}

\subsection{Pochoco SoC y procesador Espino}

El procesador Espino implementa RV32E \cite{riscv-unpriv}: el conjunto de instrucciones
base de 32 bits de RISC-V con solo 16 registros (\code{x0} a \code{x15}). No incluye
la extensión M (multiplicación y división). Además, el ALU del repositorio tiene los
desplazamientos deshabilitados para ahorrar celdas lógicas: \code{sll}, \code{srl},
\code{sra} y sus variantes con inmediato se decodifican bien, pero se ejecutan como una
suma. Estas dos restricciones determinan cómo se escribió el programa
(sección~\ref{sec:sin-mul}).

El core tiene dos etapas: búsqueda de instrucción (IF) y decodificación-ejecución (ID).
El banco de registros está hecho con bloques de RAM de la FPGA, que entregan el dato un
ciclo después de pedirlo, así que cada instrucción espera un ciclo por sus operandos.
Se midió en simulación la duración de cada tipo de instrucción, leyendo el contador de
ciclos antes y después de ella (tabla~\ref{tab:cpi}). Estos valores se usan en la
sección~\ref{sec:temporizacion} para acotar el error de la medición.

\begin{table}[H]
\centering
\caption{Ciclos de reloj por instrucción del Espino, medidos en simulación.}
\label{tab:cpi}
\begin{tabular}{lc}
\toprule
Instrucción & Ciclos \\
\midrule
ALU (\code{add}, \code{addi}, \code{xor}, \ldots), \code{lui}, \code{sw} & 2 \\
Branch no tomado & 2 \\
\code{lw} (RAM o periférico) & 3 \\
Branch tomado, \code{jal}, \code{jalr} & 3 \\
\bottomrule
\end{tabular}
\end{table}

La RAM tiene 512 palabras (2\,KB) y guarda a la vez el programa y los datos. Se
inicializa al sintetizar con el contenido de un archivo \code{.hex} mediante
\code{\$readmemh}, y el procesador comienza a ejecutar en la dirección 0. Los
periféricos se seleccionan con el bit 31 de la dirección
(tabla~\ref{tab:mapa}).

\begin{table}[H]
\centering
\caption{Mapa de memoria usado por el juego. La última fila es el registro agregado.}
\label{tab:mapa}
\begin{tabular}{lll}
\toprule
Dirección & Acceso & Contenido \\
\midrule
\code{0x0000\_0000} & lectura/escritura & RAM: programa y datos (512 palabras) \\
\code{0x8000\_0000} & escritura & displays: bits 7:4 izquierdo, bits 3:0 derecho \\
\code{0x8000\_0004} & escritura & LEDs, bits 3:0 \\
\code{0x8000\_0008} & lectura & botones, bits 3:0 (1 = apretado) \\
\code{0x8000\_000C} & lectura & \textbf{contador de ciclos, 32 bits} \\
\bottomrule
\end{tabular}
\end{table}

Cada display muestra un dígito hexadecimal: el periférico traduce los 4 bits a los
siete segmentos. Para mostrar un número decimal de dos cifras el programa escribe
$\text{decenas}\cdot 16 + \text{unidades}$; por ejemplo, para 37 escribe \code{0x37}.

\subsection{Modificaciones al RTL}

Se hicieron tres cambios al hardware, todos visibles con \code{git diff} desde el primer
commit del repositorio, que es el código original sin modificar. El tercero corrige un
defecto del SoC original que solo se manifestó en la placa
(sección~\ref{sec:placa}).

\paragraph{Contador de ciclos.} El juego necesita medir 3 segundos y tiempos de reacción.
Un bucle de espera calibrado no sirve para medir, porque su duración depende de qué
instrucciones ejecuta (tabla~\ref{tab:cpi}) y cambiaría con cualquier modificación al
programa. Se agregó a \code{rtl/pochoco\_periph.v} un registro de 32 bits que suma uno
en cada flanco de reloj y que el programa lee en el desplazamiento \code{0x0C} del
bloque de periféricos (código~\ref{lst:contador}). Es de solo lectura, comienza en cero
tras el reset y da la vuelta cada $2^{32}$ ciclos, unos 171,8 segundos.

\begin{lstlisting}[language=Verilog, caption={Contador de ciclos agregado a \code{pochoco\_periph.v}.}, label={lst:contador}]
  reg [31:0] cycle_q;

  always @(posedge clk_i or negedge rst_ni) begin
    if (!rst_ni) cycle_q <= 32'b0;
    else         cycle_q <= cycle_q + 32'd1;
  end
  ...
      case (off)
        6'd2: rdata_o <= {28'b0, btn_i}; // Buttons
        6'd3: rdata_o <= cycle_q;        // Cycle counter   <- nuevo
\end{lstlisting}

\paragraph{Módulo superior \code{game\_top}.} El archivo \code{rtl/game\_top.v} es el
que se sintetiza. Tiene los mismos puertos que la Go Board según \code{goboard.pcf} y
solo instancia el SoC con el parámetro \code{MemFile} apuntando a \code{sw/game.hex}.
El botón $i$ corresponde al LED $i$: \code{i\_Switch[i]} y \code{o\_LED[i]}.

\paragraph{Reset de encendido.} En \code{rtl/pochoco\_soc.v}, el contador que mantiene
al procesador en reset después de configurar la FPGA pasó de 4 a 12 bits: el procesador
arranca a los 4.095 ciclos (164~$\mu$s) en vez de a los 15 (0,6~$\mu$s). Con el valor
original, el juego no funcionaba en la placa aunque todas las simulaciones pasaban. La
sección~\ref{sec:placa} describe cómo se encontró la causa.

\begin{lstlisting}[language=Verilog, caption={Reset de encendido en \code{pochoco\_soc.v}, con el contador ampliado.}, label={lst:por}]
  reg  [11:0] por_cnt = 12'b0;           // antes: reg [3:0], hasta 4'hF
  wire        rst_ni;
  always @(posedge clk) begin
    if (por_cnt != 12'hFFF) por_cnt <= por_cnt + 12'd1;
  end
  assign rst_ni = (por_cnt == 12'hFFF);
\end{lstlisting}

\paragraph{Lo que no se cambió.} Se evaluó activar los desplazamientos en el ALU, pero
el SoC original ya ocupaba el 83\,\% de las celdas lógicas de la FPGA, y un
desplazador de 32 bits cuesta del orden de cien celdas. El juego se escribió sin
desplazamientos (sección~\ref{sec:sin-mul}).

\subsection{Recursos y frecuencia}

La tabla~\ref{tab:recursos} compara el SoC original con el diseño final, ambos
sintetizados con Yosys y nextpnr \cite{osscad}. El contador de ciclos agrega 96 celdas
lógicas: 32 flip-flops, la cadena de suma y la entrada nueva del multiplexor de lectura.
Los 8 bits extra del reset de encendido agregan otras 14. La frecuencia máxima que reporta
nextpnr sigue muy por encima de los 25~MHz del reloj de la Go Board, y el analizador
independiente \code{icetime} de IceStorm da un resultado equivalente. El Makefile ejecuta nextpnr con \code{--freq 25}, así que la compilación falla
si alguna vez no se alcanzara esa frecuencia.

\begin{table}[H]
\centering
\caption{Uso de la FPGA iCE40 HX1K.}
\label{tab:recursos}
\begin{tabular}{lcc}
\toprule
 & SoC original & Diseño final \\
\midrule
Celdas lógicas (de 1280) & 1063 (83\,\%) & 1173 (91\,\%) \\
Bloques de RAM (de 16) & 12 & 12 \\
Frecuencia máxima & 45,0~MHz & 42,8~MHz \\
\bottomrule
\end{tabular}
\end{table}
